\section*{FAQ}
\addcontentsline{toc}{section}{FAQ}
\markboth{FAQ}{FAQ}

\addcontentsline{toc}{subsection}{Comment entraîner une nouvelle politique ?}
\noindent \textbf{Comment entraîner une nouvelle politique ?}

\vspace{10pt}

\noindent Pour entraîner une nouvelle politique, il suffit de définir les expériences dans \texttt{experiments.yaml}, de sélectionner le mode \texttt{train} dans le fichier \texttt{main.py}, puis d'exécuter le programme avec la commande \texttt{python main.py}.

\vspace{10pt}

\noindent Il est également possible d'entraîner uniquement une partie d'un curriculum afin d'en analyser progressivement les performances. Pour cela, il suffit de modifier les paramètres \texttt{experiments} et \texttt{previous\_exp} de la fonction \texttt{run\_train()} afin de fournir la sous-liste d'expériences à entraîner ainsi que l'expérience dont les poids doivent être chargés.

\vspace{10pt}

\noindent Par exemple, le code suivant entraîne uniquement les cinq dernières expériences du curriculum :

\begin{codebox}{python}
run_train(
    experiments=train_experiments[5:],
    previous_exp=train_experiments[-6],
    trained_agent_id="agent_1",
)
\end{codebox}

\vspace{10pt}

\addcontentsline{toc}{subsection}{Quelle est la différence entre la validation et l'évaluation ?}
\noindent \textbf{Quelle est la différence entre la validation et l'évaluation ?}

\vspace{10pt}

\noindent La validation consiste à évaluer chaque politique sur le scénario sur lequel elle a été entraînée afin de vérifier les performances effectivement acquises. L'évaluation applique au contraire une même politique à plusieurs scénarios différents afin d'étudier sa capacité de généralisation.

\vspace{10pt}

\addcontentsline{toc}{subsection}{Quand utiliser \texttt{run\_demo} ou \texttt{run\_animation} ?}
\noindent \textbf{Quand utiliser \texttt{run\_demo} ou \texttt{run\_animation} ?}

\vspace{10pt}

\noindent La fonction \texttt{run\_demo} permet d'observer rapidement le comportement d'une politique au cours de son exécution, sans enregistrer de fichier. La fonction \texttt{run\_animation} enregistre quant à elle une animation (GIF et MP4), ce qui facilite l'analyse des résultats ou leur partage.

\vspace{10pt}

\addcontentsline{toc}{subsection}{À quoi sert le paramètre \texttt{trained\_agent\_id} ?}
\noindent \textbf{À quoi sert le paramètre \texttt{trained\_agent\_id} ?}

\vspace{10pt}

\noindent Ce paramètre indique quel agent a été utilisé lors de l'entraînement de la politique. Il permet ainsi de charger les poids correspondants au bon agent lorsque plusieurs agents sont présents dans un scénario.

\vspace{10pt}

\addcontentsline{toc}{subsection}{Où sont enregistrés les résultats ?}
\noindent \textbf{Où sont enregistrés les résultats ?}

\vspace{10pt}

\noindent Les poids des réseaux sont enregistrés dans le dossier \texttt{rl/SAC/weights/}. Les métriques sont sauvegardées dans le dossier \texttt{logs/} et les figures ou animations dans le dossier \texttt{figures/}. Chaque mode d'exécution possède son propre sous-dossier afin de séparer les résultats : \texttt{train/} pour le mode \texttt{train}, \texttt{validation/} pour le mode \texttt{validation}, \texttt{evaluation/} pour le mode \texttt{evaluation} et \texttt{demo/} pour le mode \texttt{animation}.

\vspace{10pt}

\addcontentsline{toc}{subsection}{Comment reprendre un entraînement ?}
\noindent \textbf{Comment reprendre un entraînement ?}

\vspace{10pt}

\noindent Lorsqu'un curriculum d'apprentissage est utilisé, les poids obtenus à la fin d'une expérience sont automatiquement réutilisés pour initialiser la suivante. Il est également possible de charger manuellement une politique existante en la renseignant dans le paramètre \texttt{previous\_exp} de la fonction \texttt{run\_train()}.

\vspace{10pt}

\noindent En revanche, il n'est pas possible de reprendre un entraînement au milieu d'une expérience. Chaque expérience doit être exécutée jusqu'à son terme (défini par l'attribut \texttt{n\_steps}) afin que les poids, les métriques et les figures correspondants soient correctement sauvegardés.

\vspace{10pt}

\addcontentsline{toc}{subsection}{Pourquoi certaines configurations semblent-elles ignorer le nombre d'obstacles ?}
\noindent \textbf{Pourquoi certaines configurations semblent-elles ignorer le nombre d'obstacles ?}

\vspace{10pt}

\noindent Certains modes de génération, comme \texttt{warehouse}, reposent sur une disposition prédéfinie des obstacles. Dans ce cas, le nombre et les dimensions des obstacles sont automatiquement déterminés afin de respecter la structure du scénario.

\vspace{10pt}

\noindent Dans les modes de génération aléatoires, comme \texttt{random}, le simulateur tente de placer le nombre d'obstacles demandé tout en respectant les contraintes géométriques, notamment la distance minimale définie par le paramètre \texttt{margin}. Si aucune configuration valide ne peut être trouvée après le nombre maximal de tentatives autorisées (\texttt{max\_attempts}), certains obstacles peuvent être omis. Il est donc possible que le scénario généré contienne moins d'obstacles que le nombre initialement demandé.

\vspace{10pt}

\addcontentsline{toc}{subsection}{Comment ajouter un nouveau scénario ?}
\noindent \textbf{Comment ajouter un nouveau scénario ?}

\vspace{10pt}

\noindent Il suffit d'ajouter une nouvelle entrée dans \texttt{environments.yaml}, puis de la référencer dans une expérience du fichier \texttt{experiments.yaml}. Cette expérience doit également être associée à une configuration d'agent et à une fonction de récompense existantes, ou à de nouvelles configurations si nécessaire. Le scénario pourra ensuite être utilisé dans tous les modes d'exécution.

\vspace{10pt}

\noindent Dans le cadre d'un curriculum d'apprentissage, il est toutefois recommandé d'utiliser la même fonction de récompense pour l'ensemble des expériences. Dans la version actuelle du simulateur, les agents sont également supposés homogènes ; il est donc conseillé de conserver une unique configuration d'agent tout au long du curriculum.
